% ICSE 2027 NIER — New Ideas and Emerging Results
% Track: NIER
% Deadline: October 23, 2026 (AoE)
% Format: IEEE conference template, 10pt, 4 pages + 1 page references
% Double-anonymous
% Compile: pdflatex icse2027-nier.tex && pdflatex icse2027-nier.tex

\documentclass[conference,10pt]{IEEEtran}

\usepackage[utf8]{inputenc}
\usepackage[T1]{fontenc}
\usepackage{amsmath,amssymb,amsthm}
\usepackage{booktabs}
\usepackage{hyperref}
\usepackage{url}
\usepackage{enumitem}
\usepackage{xcolor}

\newtheorem{theorem}{Theorem}
\newtheorem{definition}{Definition}

%% Double-anonymous: no author names
\title{Formal Cascade Propagation and Multi-Layer\\Availability Limits for In-Memory\\Infrastructure Resilience Simulation}

\author{\IEEEauthorblockN{Anonymous Author(s)}}

\begin{document}
\maketitle

%% ====================================================================
%%  ABSTRACT
%% ====================================================================
\begin{abstract}
Production chaos engineering tools inject real faults into running
systems, creating risk for regulated operators. Classical reliability
analysis (FTA, RBD) assumes component independence, which does not
hold for cloud infrastructure where correlated failures span multiple
layers simultaneously.
We present two core technical contributions implemented in an
open-source research prototype:
(1)~a cascade propagation engine formalized as a Labeled Transition
System (LTS) over a dependency graph, with proven termination,
$O(|V|{+}|E|)$ complexity, monotonicity, and bounded blast radius;
and (2)~an $N$-layer availability limit model composing categorical
ceilings via $\min$ rather than the classical series-product.
A post-hoc backtest on 54~real-world cloud incidents (2017--2024,
21~providers) achieves $F_1 = 0.87$, precision $= 1.00$, recall $= 0.81$
(severity accuracy $= 0.58$). We do not claim prospective predictive
accuracy; the backtest is illustrative. The prototype is released
under Apache~2.0 on PyPI.
\end{abstract}

%% ====================================================================
%%  1. INTRODUCTION
%% ====================================================================
\section{Introduction}\label{sec:intro}

Cascading failures in cloud infrastructure---the 2021 Meta BGP
outage~\cite{meta2021}, the 2020 AWS Kinesis cascade~\cite{aws-kinesis2020},
the 2024 CrowdStrike incident---demonstrate how a single fault
propagates across dozens of dependent services within minutes.
Two gaps remain in current approaches:

\textbf{Gap 1: Production risk.} Chaos engineering
tools~\cite{basiri2016,gremlin2020,steadybit2023,aws-fis} inject
\emph{real} faults into running systems. For regulated industries
under the EU Digital Operational Resilience Act
(DORA, EU~2022/2554), touching production to discover weaknesses
introduces unacceptable risk.

\textbf{Gap 2: Independence assumption.} Classical Reliability
Block Diagrams~\cite{trivedi2001} compose component availabilities
via the series-product $A = \prod_i A_i$, assuming statistical
independence. Cloud infrastructure violates this: a region outage
simultaneously degrades compute, network, and storage.

Concurrent work by Krasnovsky~\cite{krasnovsky2026} (ICSE-NIER~2026)
proposes in-memory Monte-Carlo simulation over service-dependency
graphs auto-discovered from distributed traces. Our work addresses
the same high-level positioning with a different technical approach:
formal LTS semantics with proven correctness properties and
multi-layer availability decomposition. The two approaches are
complementary, not competing.

\noindent\textbf{Contributions.}
\begin{enumerate}[nosep]
  \item A \textbf{cascade propagation engine} formalized as an LTS
    with proven termination, $O(|V|{+}|E|)$ complexity, monotonicity,
    causality, and bounded blast radius (\S\ref{sec:cascade}).
  \item An \textbf{$N$-layer availability limit model} using
    $\min$-composition of categorical ceilings (\S\ref{sec:avail}).
  \item A post-hoc \textbf{backtest on 54 real-world cloud incidents}
    spanning 21~providers and 7~years (\S\ref{sec:eval}).
\end{enumerate}

%% ====================================================================
%%  2. CASCADE ENGINE
%% ====================================================================
\section{Cascade Propagation Engine}\label{sec:cascade}

\begin{definition}[Infrastructure Graph]\label{def:graph}
$G = (C, E, \mathit{dep}, w, \tau, \mathit{cb})$ where $C$ is a
finite set of typed components, $E \subseteq C \times C$ directed
dependency edges, $\mathit{dep}: E \to \{\texttt{requires},
\texttt{optional}, \texttt{async}\}$, $w: E \to [0,1]$ weights,
$\tau: E \to \mathbb{R}_{\ge 0}$ latencies, and
$\mathit{cb}: E \to \{T, F\}$ circuit breaker status.
\end{definition}

\begin{definition}[Cascade Propagation Semantics (CPS)]
\label{def:cps}
An LTS $\mathcal{M} = (S, s_0, \mathit{Act}, {\longrightarrow}, F)$
where each state $s = (H, L, T, V)$ consists of:
health map $H: C \to \{\textsc{Healthy} {<} \textsc{Degraded} {<}
\textsc{Overloaded} {<} \textsc{Down}\}$,
latency map $L: C \to \mathbb{R}_{\ge 0}$,
elapsed time $T$, and visited set $V \subseteq C$
(monotonically growing).
\end{definition}

Eight transition rules govern: (R1)~fault injection,
(R2--R3)~required dependency propagation (single vs.\ multi-replica),
(R4--R5)~optional/async attenuation,
(R6)~circuit breaker trip,
(R7)~timeout cascade, and (R8)~termination.
The key correctness properties are:

\begin{theorem}[Termination]\label{thm:term}
For any finite $G$ and fault $f$, the CPS terminates. $V$ grows
monotonically and is bounded by $|C|$; for cyclic graphs, a depth
limit $D_{\max}$ and the visited set together guarantee termination
within $\min(|C|, D_{\max} \cdot \Delta)$ transitions ($\Delta$ =
max out-degree).
\end{theorem}

\begin{theorem}[Complexity]\label{thm:complex}
The CPS executes in $O(|C| + |E|)$ time: each component visited at
most once, each outgoing edge examined once.
\end{theorem}

\begin{theorem}[Monotonicity]\label{thm:mono}
Health only worsens: if $s \xrightarrow{*} s'$, then
$\forall c.\; H'(c) \ge H(c)$.
\end{theorem}

\begin{theorem}[Blast Radius Bound]\label{thm:blast}
Affected components are bounded by the reverse-reachable set:
$\mathit{affected}(c) \subseteq \{c' \mid c' \xrightarrow{*} c
\text{ in } G^{-1}\}$.
\end{theorem}

These properties distinguish the CPS from Monte-Carlo sampling
(which provides statistical but not formal guarantees) and from
static SPOF detection (which does not model propagation semantics).

\textbf{Key rules illustrated.}
Rule~2 (required dependency, single replica) captures the most
common cascade pattern: if component~$c'$ has a \texttt{requires}
dependency on failed component~$c$ with a single replica, $c'$
transitions to \textsc{Down}:
\begin{equation}
  \frac{\mathit{dep}(c', c) = \texttt{requires} \;\;
        \mathit{replicas}(c) = 1 \;\;
        H(c) = \textsc{Down}}
  {(H, L, T, V) \xrightarrow{\textit{prop}(c, c')}
   (H[c' \mapsto \textsc{Down}], L', T, V \cup \{c'\})}
\end{equation}
Rule~3 (multiple replicas) degrades rather than downs:
$\mathit{replicas}(c) > 1$ yields $H[c' \mapsto \textsc{Degraded}]$.
Rule~6 (circuit breaker) halts propagation: if
$\mathit{cb}(c', c) = \mathit{true}$ and latency exceeds the
timeout, $c'$ transitions to at most \textsc{Degraded} and no
further cascade passes through this edge. This formalization
enables reasoning about circuit breaker placement as a
graph-theoretic cut problem.

\textbf{Comparison with Krasnovsky's approach.}
Krasnovsky~\cite{krasnovsky2026} uses a ``connectivity-only
topological model'' with Monte-Carlo fail-stop sampling
(Algorithm~1 in their paper: each replica removed with probability
$p_{\text{fail}}$, BFS reachability check, 900{,}000 samples per
failure fraction). This provides \emph{statistical} resilience
estimates (MAE $\le 0.0004$ on DeathStarBench). Our CPS provides
\emph{deterministic} cascade paths with \emph{formal} guarantees
on termination and blast radius, but does not estimate resilience
distributions. The approaches address complementary questions:
``what is the expected resilience?'' vs.\ ``what is the exact
cascade path and its formal properties?''

%% ====================================================================
%%  3. N-LAYER AVAILABILITY
%% ====================================================================
\section{$N$-Layer Availability Model}\label{sec:avail}

We decompose system availability into five categorical layers:
\begin{equation}\label{eq:avail}
  A_{\text{sys}} = \min(A_{\text{L1}}, A_{\text{L2}},
  A_{\text{L3}}, A_{\text{L4}}, A_{\text{L5}})
\end{equation}
where L1 captures software availability (deploy downtime, human error,
config drift), L2 hardware (MTBF/MTTR with parallel redundancy),
L3 theoretical physics (packet loss, GC pauses), L4 operational
(incident response time, on-call coverage), and L5 external SLA
dependencies ($\prod_i a_i$).

\textbf{Why $\min$ rather than $\prod$?}
(1)~Cloud layers exhibit \emph{correlated} failure modes: a region
outage degrades L2, L3, and L5 simultaneously, violating the
independence assumption required by $\prod$.
(2)~The $\min$ operator captures a weakest-link interpretation:
observed availability over a window is dominated by whichever
category failed.
(3)~$\min$ is conservative (tighter bound) and interpretable
(the dominating layer is immediately visible).

This is a \emph{modeling choice}, not a universal theorem.
Practitioners with genuinely independent layers should use the
series-product. Within each layer, classical parallel-redundancy
formulas are retained: $A_{\text{tier}} = 1 - (1 - A_{\text{single}})^n$.

\textbf{Layer formulas.}
L1 (software): $A_{\text{L1}} = 1 - (d \cdot \bar{t}_{\text{deploy}}
/ T_{\text{month}} + p_{\text{human}} + p_{\text{drift}})$
where $d$ = deploys/month, $\bar{t}_{\text{deploy}}$ = mean deploy
downtime. L2 (hardware): per-component
$A_{\text{single}} = \text{MTBF}/(\text{MTBF}+\text{MTTR})$,
composed as $A_{\text{tier}} = 1 - (1-A_{\text{single}})^n$
with failover penalty. L4 (operational):
$A_{\text{L4}} = 1 - n_{\text{inc}} \cdot \bar{t}_{\text{resp}}
/ (8760 \cdot p_{\text{cov}})$ where $p_{\text{cov}}$ is on-call
coverage fraction. L5 (external SLA): $A_{\text{L5}} = \prod_i a_i$.

The practical insight is immediate: if
$A_{\text{L5}} = 0.999$ (three nines from external SLAs) but
$A_{\text{L2}} = 0.9999$ (four nines from hardware redundancy),
then $A_{\text{sys}} = 0.999$---the external SLA dependency
is the binding constraint, not the hardware.

\subsection{Related Work}

Fault Tree Analysis (FTA; NASA NUREG-0492, IEC~61025) and Reliability
Block Diagrams (RBD; Trivedi~\cite{trivedi2001}, IEC~61078) are
the classical analytical methods. Both operate on static trees
under component independence. HPE's SPOF patent
(US~9,280,409~B2)~\cite{hpe-spof} identifies single points of
failure in dependency graphs but does not model cascade propagation
or availability decomposition. The I$^3$ model (arXiv:2503.02890)
uses graph autoencoders for cascade prediction in urban
infrastructure---a different domain and method. Our contribution
is the specific combination of formal LTS cascade semantics,
$\min$-composition of categorical availability layers, and
open-source release for cloud infrastructure analysis.

%% ====================================================================
%%  4. EVALUATION
%% ====================================================================
\section{Evaluation}\label{sec:eval}

\textbf{Dataset.} We assembled 54~real-world cloud incidents spanning
2017--2024 from 21~providers (AWS, GCP, Azure, Meta, Cloudflare,
GitHub, Slack, Datadog, PagerDuty, DigitalOcean, Discord, Heroku,
GitLab, Atlassian, Zoom, Okta, Fastly, CrowdStrike, Dyn, OVH,
Roblox). Each incident has a
public post-mortem URL, a hand-crafted topology graph, and an
expected cascade path derived from the post-mortem narrative.

\textbf{Methodology.} For each incident, we constructed a
representative topology, injected the documented root cause, and
compared the predicted cascade path against the post-mortem.
This is a \emph{post-hoc illustrative} evaluation: topologies were
built with knowledge of the outcome. We \textbf{do not} claim
prospective prediction capability.

\textbf{Results.} Across 54~incidents: cascade path
$F_1 = 0.866$ (precision $= 1.000$, recall $= 0.810$);
severity accuracy $= 0.577$. Table~\ref{tab:provider} shows the
per-provider breakdown.
The engine processes a 50-component topology in ${<}10$\,ms on
commodity hardware.

\begin{table}[t]
\centering
\caption{Per-provider backtest summary (54 incidents).}
\label{tab:provider}
\footnotesize
\begin{tabular}{@{}lrcc@{}}
\toprule
\textbf{Provider} & \textbf{N} & \textbf{$F_1$} & \textbf{Sev.\ Acc.} \\
\midrule
AWS               & 11 & 0.846 & 0.64 \\
GCP               &  7 & 0.952 & 0.57 \\
Azure             &  6 & 0.889 & 0.57 \\
Cloudflare        &  6 & 0.797 & 0.50 \\
GitHub            &  5 & 0.880 & 0.59 \\
Other (16 prov.)  & 19 & 0.857 & 0.57 \\
\midrule
\textbf{All}      & \textbf{54} & \textbf{0.866} & \textbf{0.58} \\
\bottomrule
\end{tabular}
\end{table}

\textbf{Comparison with Krasnovsky~2026.}
Krasnovsky~\cite{krasnovsky2026} evaluates on DeathStarBench with
Monte-Carlo simulation (MAE $\le 0.0004$). The two evaluations
address different dimensions: Krasnovsky measures deviation from
observed resilience on a benchmark; our backtest measures cascade
path reproduction on real incidents. Neither constitutes prospective
validation.

\begin{table}[t]
\centering
\caption{Positioning relative to representative approaches.}
\label{tab:pos}
\footnotesize
\begin{tabular}{@{}lcccc@{}}
\toprule
& \textbf{Ours} & \textbf{K.'26} & \textbf{Chaos} & \textbf{RBD} \\
\midrule
In-memory          & \checkmark & \checkmark & \texttimes & \checkmark \\
Formal proofs      & \checkmark & \texttimes & \texttimes & partial \\
$N$-layer avail.   & \checkmark & \texttimes & \texttimes & \texttimes \\
Monte-Carlo        & \texttimes & \checkmark & \texttimes & \texttimes \\
Real injection     & \texttimes & \texttimes & \checkmark & \texttimes \\
54 incident eval.  & \checkmark & \texttimes & \texttimes & \texttimes \\
Open source        & \checkmark & \checkmark & partial    & \texttimes \\
\bottomrule
\end{tabular}
\end{table}

%% ====================================================================
%%  5. FUTURE PLANS (MANDATORY FOR NIER)
%% ====================================================================
\section{Future Plans}\label{sec:future}

This paper establishes the formal foundation and post-hoc validation
of the cascade engine and $N$-layer model. The following steps are
planned to develop this into a full-length paper:

\begin{enumerate}[nosep]
  \item \textbf{Prospective validation.} We will construct topologies
    for incidents that occurred \emph{after} the model was frozen,
    inject the root cause without knowledge of the outcome, and
    measure precision/recall on unseen incidents. This is the
    critical missing evaluation for a predictive claim.
  \item \textbf{DeathStarBench integration.} We will extract
    service-dependency graphs from DeathStarBench~\cite{gan2019}
    and compare our cascade engine's output with
    Krasnovsky's Monte-Carlo results on the same topology,
    enabling direct quantitative comparison.
  \item \textbf{Temporal layer correlation.} The current $\min$
    model treats layers as static. We plan to extend it with a
    time-varying correlation model that captures how layer
    ceilings shift during incident windows.
  \item \textbf{Formal mechanization.} We plan to mechanize the
    termination and monotonicity proofs in Coq or Lean~4 to
    eliminate the risk of proof sketch errors.
  \item \textbf{User study.} We plan a controlled study with SRE
    practitioners to evaluate whether the $N$-layer decomposition
    produces actionable insights compared to single-number
    availability metrics.
\end{enumerate}

\textbf{Prospective validation design.}
The most critical next step is prospective validation. The design
is as follows: (a)~freeze the cascade engine code at a specific
commit hash; (b)~identify 20+ cloud incidents that occurred
\emph{after} the freeze date; (c)~construct topology graphs from
public post-mortems \emph{without} consulting the affected-component
list; (d)~inject the documented root cause; (e)~measure
precision/recall against the post-mortem's affected-component list.
A blinded evaluator (not the engine developer) will construct the
topologies to prevent knowledge leakage. This protocol follows
established practice in fault localization research, where
post-hoc evaluation on known bugs is a necessary but insufficient
step, and prospective or cross-project evaluation is required for
a predictive claim.

\textbf{DeathStarBench comparison design.}
We will extract the Social Network and Hotel Reservation
service-dependency graphs from DeathStarBench~\cite{gan2019}'s
Docker Compose configurations. For each extracted graph, we will:
(a)~run our cascade engine with systematic single-fault injection
on each component; (b)~run Krasnovsky's Monte-Carlo
pipeline~\cite{krasnovsky2026} on the same graph (900{,}000
samples per failure fraction, as specified in their Algorithm~1);
(c)~compare the predicted resilience rankings and cascade paths.
This enables the first direct quantitative comparison between the
formal LTS and Monte-Carlo approaches on identical topologies.

\textbf{Formal mechanization.}
The proof sketches in this paper (Theorems~\ref{thm:term}--\ref{thm:blast})
rely on standard graph-theoretic arguments (monotone set growth,
BFS traversal bounds). Mechanizing these in Lean~4 or Coq would
strengthen the contribution by eliminating the possibility of
proof sketch errors. The Verified Software Toolchain community
has established patterns for similar graph reachability proofs
that we plan to adapt.

%% ====================================================================
%%  6. LIMITATIONS
%% ====================================================================
\section{Limitations}

The cascade engine operates on a static topology and does not
capture runtime dynamics (auto-healing, load balancer failover,
Kubernetes rescheduling). Dependency types are simplified to three
categories. The $F_1 = 0.866$ result is post-hoc and must not be
interpreted as predictive accuracy on unseen topologies. The $\min$-composition is a
modeling choice, not a universal law. The cyclic graph termination
relies on a depth limit ($D_{\max} = 20$), which is an engineering
parameter rather than a theoretically derived bound.

To be explicit, we do \emph{not} claim: (i)~predictive accuracy
on unseen topologies; (ii)~superiority over production fault
injection for empirical fault discovery; (iii)~novelty over every
prior work on cascade failure modeling; (iv)~that
$\min$-composition is universally preferable to the series-product.
What we \emph{do} claim is that the specific combination of
formal LTS cascade semantics, cross-layer $\min$-composition,
and open-source analytical implementation constitutes a useful
operating point not previously disclosed in this combined form.

\section{Discussion}

The cascade engine and $N$-layer model serve different but
complementary purposes in resilience analysis.
The cascade engine answers ``what breaks when component $X$ fails?''
with formal guarantees on the answer's completeness (bounded blast
radius) and computability (termination, linear time).
The $N$-layer model answers ``what is our availability ceiling, and
which category of failure is the binding constraint?'' with an
interpretable decomposition that makes the weakest link visible.

Together, they enable a workflow that does not exist in current
chaos engineering practice: an operator can (1)~define their
infrastructure as a dependency graph, (2)~run the cascade engine
to identify which components have the largest blast radius,
(3)~examine the $N$-layer decomposition to identify the binding
availability constraint, and (4)~prioritize resilience investments
accordingly---all without touching production.

This workflow is not a replacement for production chaos
engineering. Gremlin, Steadybit, and AWS FIS test real systems
under real load with real monitoring signals. Our engine tests
a \emph{model} of the system. The two approaches answer different
questions and should be used together where both are feasible.

%% ====================================================================
%%  REFERENCES (1 extra page allowed)
%% ====================================================================
\bibliographystyle{IEEEtran}
\begin{thebibliography}{10}

\bibitem{basiri2016}
A.~Basiri \textit{et al.},
``Chaos engineering,''
\textit{IEEE Software}, vol.~33, no.~3, pp.~35--41, 2016.

\bibitem{gremlin2020}
Gremlin, Inc.,
``Gremlin: Chaos engineering platform,'' 2020.

\bibitem{steadybit2023}
Steadybit GmbH,
``Steadybit: Chaos engineering for Kubernetes,'' 2023.

\bibitem{aws-fis}
Amazon Web Services,
``AWS Fault Injection Simulator,'' 2021.

\bibitem{trivedi2001}
K.~S.~Trivedi,
\textit{Probability and Statistics with Reliability, Queuing, and
Computer Science Applications}, 2nd~ed.
Wiley, 2001.

\bibitem{krasnovsky2026}
A.~A.~Krasnovsky,
``Model Discovery and Graph Simulation: A Lightweight Gateway to
Chaos Engineering,''
in \textit{Proc.\ ICSE-NIER}, 2026.
arXiv:2506.11176.

\bibitem{meta2021}
S.~Janardhan,
``More details about the October 4 outage,''
Meta Engineering Blog, Oct.~2021.

\bibitem{aws-kinesis2020}
Amazon Web Services,
``Summary of the Amazon Kinesis Event in the Northern Virginia
(US-EAST-1) Region,'' Nov.~2020.

\bibitem{gan2019}
Y.~Gan \textit{et al.},
``An Open-Source Benchmark Suite for Microservices,''
in \textit{Proc.\ ASPLOS}, 2019.

\bibitem{hpe-spof}
Hewlett Packard Enterprise,
``Detecting single points of failure,''
U.S. Patent 9,280,409 B2, 2016.

\end{thebibliography}

\end{document}
